% !TEX root = ../main.tex
\documentclass[../Main.tex]{subfiles}

\begin{document}
\chapter*{Forord}
Fysikk er vanskelig. Noen ganger skikkelig vanskelig. Men jeg tror likevel at de aller fleste fint kan klare å få til minimum en C og gjerne en B i et brukerkurs i ingeniørfysikk.
Men det må gjøres rett.

\section*{Forfattere}
Mens jeg selv har fungert som redaktør og hovedforfatter bak innholdet så har jeg fått fantastisk hjelp fra flinke studenter til å lage kompendiet.

\textbf{2024:} Tusen takk til Honja Kareem og Lillian Ellefsen Valla for stødig hjelp og gode innspill med å skrive inn forelesningsnotatene mine i Latex og å utfylle de der det trengtes.

\textbf{2025?:}

\section*{Foreslått framdriftsplan}

\begin{center}
\begin{tabular}{ll}
\hline
\textbf{Uke} & \textbf{Tema} \\
\hline
Uke 34 & Kapittel 1: Hva er fysikk? \\
Uke 35 & Kapittel 2: Matematiske verktøy vi trenger \\
Uke 36–37 & Kapittel 3: Kinematikk \\
Uke 38–39 & Kapittel 4: Dynamikk \\
Uke 40 & Kapittel 5: Arbeid og energi \\
Uke 41 & Kapittel 6: Statikk \\
Uke 42–43 & Kapittel 7: Fluiddynamikk \\
Uke 44–45 & Kapittel 8: Termofysikk \\
Uke 46–47 & Kapittel 9: Elektrisitet \\
\hline
\end{tabular}
\end{center}

\section{To-do-notat}
Ting som gjenstår
\begin{itemize}
    \item Utfylle teksten med eksempler og eksamensoppgaver
    \item Oppsummeringsbokser
    \item Lage videoer av forsøk / forklaringer og slenge inn QR-kode.
\end{itemize}

\end{document}